\documentclass[11pt,a4paper]{article}
\usepackage[margin=1in]{geometry}
\usepackage{amsmath,amssymb,amsthm}
\usepackage{algorithm}
\usepackage{algpseudocode}
\usepackage{booktabs}
\usepackage{hyperref}

\theoremstyle{definition}
\newtheorem{definition}{\textbf{Definition}}[section]
\newtheorem{theorem}{\textbf{Theorem}}[section]
\newtheorem{lemma}[theorem]{\textbf{Lemma}}
\newtheorem{remark}{\textbf{Remark}}[section]

\title{\textbf{Experience Replay and Rehearsal Against Catastrophic Forgetting}}
\author{\textbf{Team Scaibu} \\ \small Email: \texttt{jaydeep.9664920749@gmail.com} \\ \small \textbf{Architecture and Core Engine Team}}
\date{\textbf{October 2026}}

\begin{document}
\maketitle

\section{Definitions and Cognitive Mental Models}

\subsection{Formal Mathematical Definition}
\textbf{Definition (Experience Replay and Reservoir Memory Buffering):}
Experience replay (rehearsal) in continual learning maintains a bounded episodic memory buffer $\mathcal{M}$ of historical training experiences ($|\mathcal{M}| \le K_{\text{capacity}}$) sampled from earlier task distributions $\{\mathcal{D}_1, \dots, \mathcal{D}_{t-1}\}$. When training on a stream of incoming data from new task $\mathcal{D}_t$, mini-batches are synthesized by interleaving incoming new data samples $\mathcal{B}_{\text{new}}$ with rehearsal samples $\mathcal{B}_{\text{replay}} \sim \mathcal{M}$:
{\boldmath
\begin{equation}
\mathcal{B}_{\text{blended}} = \mathcal{B}_{\text{new}} \cup \mathcal{B}_{\text{replay}}, \quad \frac{|\mathcal{B}_{\text{replay}}|}{|\mathcal{B}_{\text{blended}}|} = \rho_{\text{replay}}
\end{equation}
}
where $\rho_{\text{replay}} \in [0, 1]$ is the rehearsal mixing ratio.

To guarantee that the memory buffer maintains a uniform, unbiased sample of all seen historical experiences without knowing the total stream length $N_{\text{total}}$ in advance, \textbf{Reservoir Sampling} (Vitter, 1985) is applied. Given buffer capacity $K$, the first $K$ stream items fill the buffer. For each subsequent incoming item index $n > K$, the item is accepted into the buffer with probability:
{\boldmath
\begin{equation}
P(\text{accept item } n) = \frac{K}{n}
\end{equation}
}
and replaces an existing item chosen uniformly at random from indices $\{1, \dots, K\}$. This guarantees that at any point in time, every item seen so far has an exact equal probability $\frac{K}{n}$ of residing in the buffer.

\subsection{Expanded Non-Technical Definition (Intuition for Anyone)}
\begin{enumerate}
    \item \textbf{Spaced Repetition Flashcards}: When learning a third foreign language, you don't throw away your vocabulary flashcards for the first two languages. You review 5 old flashcards for every 10 new words you learn! That small reminder prevents earlier memories from fading away.
    \item \textbf{Generative Replay When Data Cannot Be Stored}: Under strict privacy laws (GDPR/HIPAA), raw customer data cannot be permanently stored in a buffer. In \textbf{Generative Replay}, an internal generative model produces synthetic practice examples mimicking the old distribution without storing private user records.
    \item \textbf{Simple, Universal, and Robust}: Unlike complex regularizers that rely on Gaussian approximations of parameter loss surfaces, replay works directly on the data distribution, making it the most resilient defense against forgetting in deep neural networks.
\end{enumerate}

\subsection{Child-Friendly Mental Model (Physical Metaphor)}
Imagine a photo album with only 100 pages ($K=100$):
\begin{itemize}
    \item \textbf{FIFO (Throwing away early memories)}: Every time you take a new photo, you rip out the oldest childhood photo and throw it away. After 10 years, you have zero childhood memories left!
    \item \textbf{Reservoir Replay}: Whenever you take a new photo, you roll a special dice. The dice decides whether to replace an old photo, ensuring your 100-page album contains an even mix of childhood, teenage, and adult photos forever!
\end{itemize}

\section{Core Mathematical Equations and Definitions}

\subsection{Blended Optimization Objective}
{\boldmath
\begin{equation}
\mathcal{L}_{\text{blended}}(\theta) = (1 - \rho) \, \mathbb{E}_{x \sim \mathcal{B}_{\text{new}}} [\mathcal{L}(x; \theta)] + \rho \, \mathbb{E}_{x \sim \mathcal{M}} [\mathcal{L}(x; \theta)]
\end{equation}
}
\begin{itemize}
    \item \textbf{Formal Definition}: Weighted loss evaluated over blended union of current streaming data and historical buffer.
    \item \textbf{What It Means}: Forces model weights to minimize current task error while remaining anchored to old task loss basins.
    \item \textbf{Why It Is Important}: Directly anchors gradient updates across the union of task distributions.
\end{itemize}

\subsection{Reservoir Replacement Probability}
{\boldmath
\begin{equation}
P(X_n \in \mathcal{M}) = \frac{K}{n}, \quad n \ge K
\end{equation}
}
\begin{itemize}
    \item \textbf{Formal Definition}: Asymptotic inclusion probability for streaming item $n$ in buffer of capacity $K$.
    \item \textbf{What It Means}: Guarantees uniform memory representation across time.
    \item \textbf{Why It Is Important}: Eliminates recency bias and prevents memory saturation by newer tasks.
\end{itemize}

\section{Rigorous Mathematical Analysis Checklist}

\begin{itemize}
    \item \textbf{Notation}: $\mathcal{M}$ (memory buffer), $K$ (capacity limit), $\mathcal{B}_{\text{new}}$ (new batch), $\rho \in [0, 1]$ (replay ratio), $n$ (stream count).
    \item \textbf{Epistemic Status}: Vitter (ACM 1985); Rolnick et al. (NeurIPS 2019 - CLEAR). Gold standard empirical baseline for continual learning.
    \item \textbf{Dependency Graph}: Stream items $\to \text{Sample Replay from } \mathcal{M} \to \text{Blend Batch} \to \text{Update Weights} \to \text{Reservoir Update } \mathcal{M}$.
    \item \textbf{Dimensions}: Buffer capacity $K$ (e.g. $K \in [500, 10000]$ samples).
    \item \textbf{Regimes}: Continual domain instruction tuning, streaming financial market models, long-lived autonomous agents.
    \item \textbf{Invariants}: Buffer size strictly satisfies $|\mathcal{M}| \le K$ at all timesteps.
    \item \textbf{Isomorphisms}: Isomorphic to importance-weighted coreset construction and experience replay in reinforcement learning (DQN).
\end{itemize}

\section{Theoretical Theorems and Formal Proofs}

\begin{theorem}[Uniform Distribution Invariance of Reservoir Sampling]
Let a stream of $n$ items $x_1, x_2, \dots, x_n$ ($n \ge K$) be processed into a buffer of capacity $K$ using Algorithm $R$. For any step $n$ and any item $i \in \{1, \dots, n\}$:
{\boldmath
\begin{equation}
P(x_i \in \mathcal{M}_n) = \frac{K}{n}
\end{equation}
}
\end{theorem}

\begin{proof}
We prove this by induction on the sequence length $n$.

\textbf{Base Case ($n = K$):}
The first $K$ items are unconditionally placed into the buffer $\mathcal{M}_K = \{x_1, \dots, x_K\}$.
Thus, for every $i \in \{1, \dots, K\}$, $P(x_i \in \mathcal{M}_K) = 1 = \frac{K}{K}$, satisfying the base case.

\textbf{Inductive Hypothesis:}
Assume the theorem holds for $n - 1 \ge K$, meaning that for all $i \in \{1, \dots, n - 1\}$:
{\boldmath
\begin{equation}
P(x_i \in \mathcal{M}_{n-1}) = \frac{K}{n - 1}
\end{equation}
}
\textbf{Inductive Step for step $n$:}
Consider the $n$-th item $x_n$:
By definition of the algorithm, item $x_n$ is added to the buffer with probability $\frac{K}{n}$.
Thus:
{\boldmath
\begin{equation}
P(x_n \in \mathcal{M}_n) = \frac{K}{n}
\end{equation}
}
Now consider any previously seen item $x_i$ where $1 \le i \le n - 1$.
Item $x_i$ resides in $\mathcal{M}_n$ if and only if $x_i$ was in $\mathcal{M}_{n-1}$ AND $x_i$ was not evicted by $x_n$.
Item $x_i$ is evicted if $x_n$ is accepted (probability $\frac{K}{n}$) and $x_i$'s slot is chosen for replacement (probability $\frac{1}{K}$).
The probability of eviction given that $x_i \in \mathcal{M}_{n-1}$ is:
{\boldmath
\begin{equation}
P(\text{evicted} \mid x_i \in \mathcal{M}_{n-1}) = \frac{K}{n} \cdot \frac{1}{K} = \frac{1}{n}
\end{equation}
}
The probability that $x_i$ survives given $x_i \in \mathcal{M}_{n-1}$ is:
{\boldmath
\begin{equation}
P(\text{survives} \mid x_i \in \mathcal{M}_{n-1}) = 1 - \frac{1}{n} = \frac{n - 1}{n}
\end{equation}
}
By the law of total probability:
{\boldmath
\begin{equation}
P(x_i \in \mathcal{M}_n) = P(x_i \in \mathcal{M}_{n-1}) \cdot P(\text{survives} \mid x_i \in \mathcal{M}_{n-1}) = \frac{K}{n - 1} \cdot \frac{n - 1}{n} = \frac{K}{n}
\end{equation}
}
This completes the induction. Every item seen has identical probability $\frac{K}{n}$ of remaining in memory.
\end{proof}

\begin{theorem}[Pareto-Optimal Bound on Rehearsal Sample Complexity]
Let task distributions $\mathcal{D}_1, \dots, \mathcal{D}_T$ have pairwise distribution shifts bounded by Total Variation distance $\operatorname{TV}(\mathcal{D}_i, \mathcal{D}_j) \le \delta$. To bound catastrophic degradation below $\epsilon$ with probability $1 - \delta_0$, the required replay buffer capacity scales sublinearly as:
{\boldmath
\begin{equation}
K \ge \mathcal{O}\left( \frac{d_{\text{VC}}}{\epsilon^2} \log \frac{1}{\delta_0} \right)
\end{equation}
}
where $d_{\text{VC}}$ is the Vapnik-Chervonenkis dimension of the hypothesis class.
\end{theorem}

\begin{proof}
By uniform convergence theory in statistical learning theory:
1. The empirical risk on buffer $\mathcal{M}$ approximates the true risk over past tasks $\mathcal{R}_{\text{past}}(f) = \frac{1}{T-1} \sum_{t=1}^{T-1} \mathbb{E}_{x \sim \mathcal{D}_t}[\ell(x; f)]$ up to generalization gap:
{\boldmath
\begin{equation}
\sup_{f \in \mathcal{F}} \left| \mathcal{R}_{\mathcal{M}}(f) - \mathcal{R}_{\text{past}}(f) \right| \le \mathcal{O}\left( \sqrt{\frac{d_{\text{VC}} \log(K) + \log(1/\delta_0)}{K}} \right)
\end{equation}
}
2. Setting the right-hand bound $\le \epsilon$ and solving for $K$ yields:
{\boldmath
\begin{equation}
K \ge \mathcal{O}\left( \frac{d_{\text{VC}}}{\epsilon^2} \log \frac{1}{\delta_0} \right)
\end{equation}
}
This establishes that a fixed, constant-sized buffer is mathematically sufficient to bound catastrophic forgetting arbitrarily close to zero regardless of training duration.
\end{proof}

\section{Foundational Architectural Questions and Epistemic Meta-Questions}

\subsection{Architectural Question 1: Buffer Selection Policies}
\textbf{Question:} Why is reservoir sampling preferred over simple FIFO queue eviction?\\
\textbf{Answer:} FIFO queues discard the oldest data. When sequentially training across Tasks 1, 2, 3, and 4, FIFO completely purges Task 1 by the time Task 4 arrives, causing catastrophic forgetting of earliest skills. Reservoir sampling preserves a balanced cross-section of all historical tasks throughout training.

\textbf{Meta-Question:} How can replay samples be selected for maximum informational value?\\
\textbf{Answer:} In \textbf{Gradient Episodic Memory (GEM)} and \textbf{Maximal Interfering Retrieval (MIR)}, samples are selected not purely at random, but by prioritizing past examples whose losses are predicted to increase the most under the candidate parameter update.

\subsection{Architectural Question 2: Replay in Generative LLMs}
\textbf{Question:} How is replay applied when fine-tuning LLMs on enterprise domains (e.g. legal or medical documents)?\\
\textbf{Answer:} Domain fine-tuning of LLMs without replay causes rapid degradation of general chat abilities and formatting skills. Standard practice mixes $10\%\text{--}20\%$ of general instruction tuning data (e.g. Alpaca, UltraChat, or pretraining web text) into every domain training batch.

\textbf{Meta-Question:} What are the regulatory constraints of raw data rehearsal?\\
\textbf{Answer:} Rehearsal buffers storing personal data violate GDPR ``right to be forgotten'' regulations. When a user requests deletion, their entries must be scrubbed from the buffer, or generative pseudo-rehearsal must be utilized instead of raw sample caching.

\section{Algorithmic Implementation Specification}

\begin{algorithm}[H]
\caption{Replay Buffer Reservoir Update and Batch Blending}
\begin{algorithmic}[1]
\Require New batch $\mathcal{B}_{\text{new}}$, memory buffer $\mathcal{M}$, replay ratio $\rho$, capacity $K$, seen count $N_{\text{seen}}$
\Ensure Blended batch $\mathcal{B}_{\text{blended}}$, updated buffer $\mathcal{M}'$, replay count $k_{\text{rep}}$, new count $k_{\text{new}}$
\State $k_{\text{new}} \gets |\mathcal{B}_{\text{new}}|$
\If{$|\mathcal{M}| > 0$ \textbf{and} $\rho > 0.0$}
    \State $k_{\text{rep}} \gets \min(|\mathcal{M}|, \lceil k_{\text{new}} \cdot \rho / (1.0 - \rho) \rceil)$
    \State Sample $k_{\text{rep}}$ items from $\mathcal{M} \to \mathcal{B}_{\text{replay}}$
\Else
    \State $\mathcal{B}_{\text{replay}} \gets \emptyset, \quad k_{\text{rep}} \gets 0$
\EndIf
\State $\mathcal{B}_{\text{blended}} \gets \mathcal{B}_{\text{replay}} \cup \mathcal{B}_{\text{new}}$
\State $\mathcal{M}' \gets \operatorname{Copy}(\mathcal{M})$
\For{\textbf{each} sample $x \in \mathcal{B}_{\text{new}}$}
    \State $N_{\text{seen}} \gets N_{\text{seen}} + 1$
    \If{$|\mathcal{M}'| < K$}
        \State Append $x$ to $\mathcal{M}'$
    \Else
        \State Draw replacement index $j \gets \operatorname{Hash}(N_{\text{seen}}) \pmod{N_{\text{seen}}}$
        \If{$j < K$}
            \State $\mathcal{M}'[j] \gets x$
        \EndIf
    \EndIf
\EndFor
\State \Return $(\mathcal{B}_{\text{blended}}, \mathcal{M}', k_{\text{rep}}, k_{\text{new}})$
\end{algorithmic}
\end{algorithm}

\section{Big-$\Theta$ Complexity Analysis}

\begin{itemize}
    \item \textbf{Replay Batch Sampling Time}: $\Theta(k_{\text{rep}})$ operations to fetch rehearsal samples.
    \item \textbf{Reservoir Maintenance Time}: $\Theta(k_{\text{new}})$ hash and assignment operations.
    \item \textbf{Batch Synthesis Time}: $\Theta(k_{\text{new}} + k_{\text{rep}})$ pointer aggregations.
    \item \textbf{Total Time Complexity}: $\Theta(|\mathcal{B}_{\text{blended}}|)$ operations per training iteration.
    \item \textbf{Buffer Space Complexity}: $\Theta(K)$ memory slots storing rehearsal records.
\end{itemize}

\section{Single Canonical Repository Reference}

The complete deterministic scalar implementation, verification suite, and unit test benchmarks are published at:
\begin{center}
\url{https://github.com/Chief-Strategist-J/llm-obs-infra}
\end{center}

\end{document}
